\section{Conclusão}

As evidências examinadas sustentam que a taxa de resolução, isoladamente, não basta para comparar agentes de programação sob restrições de engenharia. O consumo de recursos, o harness e o procedimento de entrega afetam a interpretação do desempenho; exposição prévia às tarefas, verificadores incompletos e seleção de subconjuntos limitam o alcance das conclusões. Esses fatores fundamentam a avaliação conjunta de qualidade operacional, custo e latência.

Os estudos de arquitetura fundamentam investigar conjuntamente modelos, contexto e processo. Composição fixa por atividade, roteamento condicional e seleção de informação constituem decisões distintas; a vantagem de separá-las ou combiná-las depende da tarefa, da capacidade do modelo e da confiabilidade da verificação. A evidência disponível motiva essas intervenções, sem demonstrar que aumentar o número de agentes, fornecer mais contexto ou adicionar etapas produza benefício uniforme.

O protocolo estabelece como avaliar uma abordagem frente ao modelo no harness padrão e a controles de menor custo, economia de tokens no modelo forte e escalonamento simples. Contrastes de composição, contexto e processos fixos serão associados aos mecanismos efetivamente adotados. O benefício será examinado tanto pela preservação da resolução com economia quanto por uma redução substancial de custo que justifique uma perda delimitada de resolução, sob restrição de tempo. Piloto e desenvolvimento antecederão o congelamento das configurações e a avaliação reservada. Composição e roteamento de modelos, engenharia de contexto, ativação de etapas, revisão e paralelização permanecem hipóteses, sem superioridade experimental demonstrada para a abordagem.
